\documentclass[a4paper,openright,10pt]{article}
\usepackage{amsmath}
\usepackage{amssymb}
\usepackage{graphicx}
\usepackage{fancyhdr}
\pagestyle{fancy}
\fancyhf{}
\rfoot{P\'agina \thepage \hspace{1pt} de 2}
\everymath{\displaystyle}

\begin{document}

\title{Probabilidad y Estadística}
\author{Prof. Mtro. Luis Jos\'{e} Yudico Anaya}
\date{15-10-2024}
\maketitle

\begin{center}
\textbf{INGENIERÍA EN NANOTECNOLOGÍA}\\
\vspace{.5cm}
\textbf{2024-2025-I}\\
\vspace{.5cm}
\textbf{EXAMEN SOBRE PROBABILIDAD BÁSICA}
\end{center}


\title{Ejercicios de Probabilidad en Nanotecnología}
\author{Para alumnos de tercer semestre}
\date{}
\maketitle

\section*{Ejercicio 1: Espacios Muestrales}
En un experimento de síntesis de nanopartículas, se producen 100 nanopartículas de diferentes tamaños. Si 30 de ellas tienen un tamaño menor a 10 nm, define el espacio muestral y determina la probabilidad de seleccionar una nanopartícula menor a 10 nm.

\section*{Ejercicio 2: Eventos}
Un investigador estudia la efectividad de un tratamiento en 50 muestras de material a nanoescala. Si 20 muestras responden positivamente al tratamiento, ¿cuál es la probabilidad de que, al seleccionar al azar dos muestras, ambas respondan positivamente?

\section*{Ejercicio 3: Experimentos con Resultados Igualmente Probables}
En un estudio sobre la toxicidad de nanopartículas, se seleccionan aleatoriamente 4 tipos de nanopartículas. Si cada tipo tiene una probabilidad de toxicidad del 15\%, ¿cuál es la probabilidad de que al menos una de las nanopartículas sea tóxica?

\section*{Ejercicio 4: Funciones de Masa de Probabilidad}
Un nuevo tipo de nanopartícula tiene las siguientes probabilidades de causar efectos adversos en un experimento: \( P(efecto \, adverso) = 0.1 \), \( P(efecto \, neutro) = 0.8 \), \( P(effecto \, positivo) = 0.1 \). ¿Cuál es la probabilidad de que el efecto sea neutro o positivo?

\section*{Ejercicio 5: Probabilidad Condicional}
En un estudio de nanopartículas, el 60\% de ellas son metálicas y el 70\% de las metálicas tienen propiedades conductoras. Si seleccionas al azar una nanopartícula que tiene propiedades conductoras, ¿cuál es la probabilidad de que sea metálica?

\section*{Ejercicio 6: Teorema de Bayes}
Un laboratorio desarrolla una prueba para detectar nanopartículas contaminantes. La prueba tiene una tasa de detección del 95\% para nanopartículas contaminantes y un 5\% de falsos positivos. Si el 2\% de las muestras contiene nanopartículas contaminantes, ¿cuál es la probabilidad de que una muestra positiva realmente contenga nanopartículas contaminantes?

\section*{Ejercicio 7: Eventos Independientes}
Un proceso de fabricación de nanomateriales tiene un 80\% de probabilidad de éxito en la creación de una muestra y un 90\% de probabilidad de que dicha muestra tenga las propiedades deseadas. ¿Cuál es la probabilidad de que el proceso sea exitoso y la muestra tenga las propiedades deseadas, suponiendo que los eventos son independientes?

\section*{Ejercicio 8: Regla de Multiplicación}
En un estudio sobre la eficacia de nanopartículas en la administración de fármacos, la probabilidad de que una nanopartícula se adhiera a una célula es del 75\% y la probabilidad de que entregue el fármaco es del 85\%. ¿Cuál es la probabilidad de que la nanopartícula se adhiera a una célula y entregue el fármaco?

\section*{Ejercicio 9: Aplicaciones de la Probabilidad}
En una producción de 1000 nanopartículas, se estima que el 5\% tiene defectos. Si se seleccionan 20 nanopartículas al azar, ¿cuál es la probabilidad de que al menos una de ellas tenga un defecto?

\section*{Ejercicio 10: Combinatoria y Probabilidad}
Un equipo de investigación tiene 4 tipos de nanopartículas para estudiar. Si seleccionan 3 de ellas para un experimento, ¿cuál es la probabilidad de que al menos 2 de los tipos seleccionados sean del mismo grupo (por ejemplo, metálicas o no metálicas)?

\end{document}
